\centerline{\bf \Huge Домашнее задание}
\centerline{\bf \Huge EZ}

\begin{flushright}
        \scriptsize{Хоть задач и 7, но если решили больше 5, то сверхбаллы можно потратить, чтобы закрыть старые штрафные задачи!!!}
\end{flushright}

\problem{1}
На олимпиаду по математике у каждого школьника пришло не более 5 одноклассников. Докажите, что всех ребят можно рассадить по трем кабинетам так, чтобы ни в одном кабинете не оказались три человека из одного класса.

\problem{2}
Имеются булочки 4 видов - по 6 булочек каждого вида. Их разделили поровну между шестью детьми. Какое наименьшее количество детей можно выбрать так, чтобы у них гарантированно нашлись булочки всех видов?


\problem{3}
Каждое из 8 различных натуральных чисел не превосходит 14. Докажите, что среди их попарных разностей есть хотя бы 3 одинаковые.

\problem{4}
В лесу растет миллион елок. Известно, что на каждой из них не более 600000 иголок. Докажите, что в лесу найдутся две елки с одинаковым числом иголок.

\problem{5}
В школе 1500 учащихся. Найдется ли такой день в году, в который отмечают свой день рождения не меньше, чем 5 учеников?

\problem{6}
Пусть в классе 41 человек. Некто сделал в контрольной работе больше всех ошибок - 13. Докажите, что найдутся 4 учащихся, сделавших одинаковое число ошибок.

\problem{7}
Все точки плоскости покрашены в один из двух цветов. Докажите, что найдутся две точки на расстоянии 1 метр друг от друга, окрашенные одинаково.